\section{Conclusion}

We introduced Anchor Transfer Learning for drug--target affinity prediction, a framework that conditions each prediction on an anchor protein already known to bind the same drug. Rather than scoring protein--drug pairs in isolation, the method turns affinity prediction into a comparison against experimentally grounded binding evidence. The formulation requires no structural information and applies to any target for which binder annotations exist.

The central empirical finding is a generalization gap: models that rank highest within a single benchmark are not the models that transfer best. DeepDTA and ESM-DTA collapse to random performance on Davis kinases (per-protein CI 0.520 and 0.501 respectively), despite achieving strong same-dataset metrics on DTC. Anchor transfer reverses this pattern: V2-650M achieves per-protein CI of 0.642 and AUROC of 0.669 on Davis. Crucially, anchors are retrieved at test time from the training set by Tanimoto chemical similarity after excluding all canonical chemical duplicates (83.8\% of Davis drugs were present in DTC under different SMILES strings), ensuring genuine zero molecular overlap. Performance stratification by retrieval similarity confirms that the model generalizes even to chemically dissimilar compounds (Tanimoto~$< 0.6$: CI 0.656).

A cosine similarity ablation confirms that the model learns more than protein similarity: 74\% of V2-650M's predictive variance is orthogonal to embedding cosine distance, and the full model outperforms cosine similarity on 82\% of Davis proteins. The three-way interaction architecture captures drug-specific binding patterns that raw protein similarity cannot encode.

ESM-DTA serves as a critical negative control: it is the strongest same-dataset model (AUROC 0.910 on DTC) yet collapses to random on Davis (per-protein AUROC 0.503). This demonstrates that stronger protein embeddings alone do not solve the cross-dataset transfer problem---the anchor mechanism provides the missing ingredient by conditioning predictions on experimentally grounded binding evidence.

The method is designed for the common drug discovery setting where at least one strong binder of the query drug is already known---lead optimization, scaffold hopping, selectivity profiling, and repurposing all satisfy this condition. Several directions remain open: learned anchor retrieval to improve on the current Tanimoto nearest-neighbor strategy, multi-anchor ensembles for calibrated uncertainty, deeper cross-attention architectures for richer three-way interactions, and evaluation on structurally diverse protein families beyond kinases.
